\chapter{Base Change and Generated Comparisons}\label{chap:5}

\section*{Overview of the Chapter}

The question of this chapter is \textbf{whether the operation that carries structure and the
operation that changes the input under consideration can be interchanged}.
Chapter~\ref{chap:4} treated transport along a single change and the coherence of successive
transports. In this chapter we construct the two routes formed from these two kinds of
operations and study the comparison between them.

For example, suppose that a monolith with order processing and inventory management is split
into an order service and an inventory service. In a design review we first want to examine
the part related to order cancellation. That part contains the operations of cancellation
request, release of an inventory reservation, and cancellation completion, together with the
condition that if a cancellation is complete then the corresponding inventory reservation has
been released. Under the same splitting policy, we compare two routes that yield the design of
this part.

\begin{xltabular}{\linewidth}{@{}LLL@{}}
\toprule
Procedure & First step & Next step \\
\midrule
\endhead
Split the whole, then select the scope & Split the whole monolith into the order and inventory
services & Extract the part related to order cancellation from the design after the split \\
\addlinespace[3pt]
Select the scope, then split & Extract the part related to order cancellation from the
original design & Apply the same splitting policy to that part \\
\bottomrule
\end{xltabular}

Does the splitting proposal that the person in charge drew up by focusing on the order-cancellation
part remain consistent with the design obtained after splitting the whole? For example,
consider a design in which, after the split, the order service asks the inventory service to
release the reservation and proceeds to cancellation completion only after confirming the
release. Between the designs obtained by the two routes, we want to verify that this
interaction and the condition above correspond. Extracting only the processing placed in the
order service leaves the release of inventory reservations outside the scope, and we lose the
dependencies needed to check the condition.

We therefore look for the inputs on the original that correspond to the scope selected on the
target. In this example, this means tracing the operations and dependencies needed for the
cancellation condition back to the design before the split. In this chapter we construct that
correspondence as pairs of inputs and pull the structure of objects, operations, and Laws back
onto it. This is base change. Moreover, from the universal properties of transport and
pullback, we build a map that compares the results of the two routes. We call this map the
canonical comparison.

\begin{xltabular}{\linewidth}{@{}LLL@{}}
\toprule
Question & Construction & Outcome \\
\midrule
\endhead
Can two inputs be matched under the same condition? & Pairs of sources with a common image &
Fiber products of extraction doctrines \\
\addlinespace[3pt]
Can the structure on the target be read on the original? & Backward reindexing of Atoms &
Cartesian pullback and its equivalence with transport \\
\addlinespace[3pt]
Can the order of transport and pullback be exchanged? & Squares generated from finite codes,
and mates & A canonical invertible comparison \\
\addlinespace[3pt]
Can several comparisons and diagnoses also be carried? & Commutative base diagrams and the
correspondence of end groups & Preservation and reflection of coherence, vanishing of defects,
and reselection orbits \\
\addlinespace[3pt]
Can structure still be carried back even along a change that enlarges the extracted content? & Reflection of extraction at
inputs that actually carry structure & A necessary and sufficient condition for backward
transport along refinements \\
\addlinespace[3pt]
What is compared when geometry and normalization are included? & Two routes of full geometries
built from a finitely presented square and the same comparison diagram & The generated
comparison analyzed in Chapter~\ref{chap:6} \\
\bottomrule
\end{xltabular}

For the exact changes of the first half, the Atoms extracted from the selected sources agree,
including the change of names. For the refinements of the second half, only forward
preservation of the extracted content is required. This difference determines which structure
can be carried backwards. When treating the split of a monolith as well, we specify whether we
read the business operations and conditions or also communication, failures, and retries, and
we check the preservation conditions for extraction and structure in that reading.

\section{Constructing corresponding inputs by fiber products}\label{sec:5.1}

An extraction doctrine is a reading that normalizes a source and then extracts Atoms. The
category $B$ of Chapter~\ref{chap:1} has objects that in addition select one source. We first
carry out the construction in the category of doctrines with the selected source forgotten,
and then restore the selection.

\begin{construction}[Fiber product of exact doctrines]\label{cons:5.1}

Over a common Atom set $\mathrm{At}$, take two exact morphisms.
\begin{equation*}
D_1\xrightarrow{\sigma_1}D_0\xleftarrow{\sigma_2}D_2,
\qquad \sigma_i=(f_i,e_i).
\tag{5.1}\label{eq:5.1}
\end{equation*}
Here $f_i$ is the source map and $e_i$ is the bijection of Atoms. We define the sources and
the normalization of the new doctrine $D_{12}$ by
\begin{equation*}
\begin{aligned}
\mathrm{Src}_{12}
 &=\{(s_1,s_2)\mid f_1(s_1)=f_2(s_2)\},\\
N_{12}(s_1,s_2)&=(N_1s_1,N_2s_2)
\end{aligned}
\tag{5.2}\label{eq:5.2}
\end{equation*}
The normalization preserves this set because $f_iN_i=N_0f_i$. The choices of vocabulary,
semantics, and resolution are taken from $D_1$, and the predicates that read a source receive
the first component. Hence
\begin{equation*}
\mathrm{Extracts}_{12}((s_1,s_2),a)
\Longleftrightarrow \mathrm{Extracts}_1(s_1,a).
\tag{5.3}\label{eq:5.3}
\end{equation*}
The two projections are
\begin{equation*}
\pi_1=(\mathrm{pr}_1,\mathrm{id}_{\mathrm{At}}),
\qquad
\pi_2=(\mathrm{pr}_2,e_2^{-1}e_1)
\tag{5.4}\label{eq:5.4}
\end{equation*}
The Atom map of the second projection carries the first names to the names of the common base
and returns from there to the second names. Exactness follows from
\begin{equation*}
\begin{aligned}
\mathrm{Extracts}_1(s_1,a)
&\Longleftrightarrow\mathrm{Extracts}_0(f_1(s_1),e_1(a))\\
&\Longleftrightarrow\mathrm{Extracts}_2(s_2,e_2^{-1}e_1(a))
\end{aligned}
\tag{5.5}\label{eq:5.5}
\end{equation*}
In both the source and the Atom components we have $\sigma_1\pi_1=\sigma_2\pi_2$.

\end{construction}

\begin{proposition}[Universal property with respect to all cones]\label{prop:5.2}

If exact morphisms $h_i:T\to D_i$ satisfy $\sigma_1h_1=\sigma_2h_2$, then there is a unique
exact morphism $h:T\to D_{12}$ with $\pi_i h=h_i$. Hence $D_{12}$ is the fiber product
of~\eqref{eq:5.1}.

\end{proposition}

\begin{proof}

Write $h_i=(g_i,d_i)$, and let the source map of $h$ be $t\mapsto(g_1(t),g_2(t))$ and its
Atom map be $d_1$. The pair of sources belongs to~\eqref{eq:5.2} by the commutativity of the
cone. Commutation with the normalizations holds componentwise, and the equivalence of
extraction is the exactness of $h_1$. For the second projection, from $e_1d_1=e_2d_2$ we
obtain $e_2^{-1}e_1d_1=d_2$. The equations with the two projections determine the source map
uniquely, and the equation with the first projection determines the Atom map uniquely.

\end{proof}

When the selected sources $s_i$ satisfy $f_1(s_1)=f_2(s_2)$, we select $(s_1,s_2)$ in
$D_{12}$. If the two morphisms of a cone with specified selected sources (a pointed cone)
preserve the selected sources, then so does the universal morphism above. In this way the
same fiber product is obtained in $B$.

\begin{example}[Inputs corresponding to a scope on the target]\label{ex:5.3}

We check Construction~\ref{cons:5.1} in a case where the sources form small finite sets. Let
the sources of the original be $\{a,b,c\}$ and the sources of the target be $\{0,1\}$, with
the correspondence $f(a)=f(b)=0$ and $f(c)=1$. The fiber product along the inclusion
$j:\{0\}\to\{0,1\}$, which selects the object $\{0\}$ after the change, is
\begin{equation*}
\{a,b,c\}\times_{\{0,1\}}\{0\}
=\{(a,0),(b,0)\}.
\tag{5.6}\label{eq:5.6}
\end{equation*}
If every source extracts the same finite Atom set and the normalizations and the Atom maps
are identities, this becomes a concrete instance of Construction~\ref{cons:5.1}. Even though
$a,b$ are sent to the same source of the target, both remain as corresponding inputs of the
original. When applying this to the split of a monolith, we take design documents and code as
sources, specify the reading used for the design change and the selection of the scope, and
check~\eqref{eq:5.5}.

\end{example}

\section{Pulling back cores from the target}\label{sec:5.2}

The transport $\sigma_!$ of Chapter~\ref{chap:4} started from a core on the original. The
pullback $\sigma^*$ starts from a core on the target. Both correspond to the same base
morphism $\sigma:b\to b'$, but the direction in which the object is constructed is reversed.

\begin{definition}[Strongly cartesian morphism]\label{def:5.4}

Take a functor $r:E\to B$ and a morphism $\kappa:Y\to Q$ with $r(\kappa)=\sigma:b\to b'$. We
call $\kappa$ a strongly cartesian morphism when, for every object $R$, base morphism
$\tau:r(R)\to b$, and morphism $h:R\to Q$ with $r(h)=\sigma\tau$,
\begin{equation*}
\exists!\,k:R\longrightarrow Y,\qquad
r(k)=\tau,\quad \kappa k=h
\tag{5.7}\label{eq:5.7}
\end{equation*}
This is the condition that a morphism into $Q$ splits uniquely into a final step
corresponding to $\sigma$ on the base and a morphism up to that point
(\cite[\href{https://stacks.math.columbia.edu/tag/02XJ}{Tag 02XJ, Definition 4.33.1}]{Stacks}).

\end{definition}

\begin{construction}[Backward reindexing]\label{cons:5.5}

Let $\sigma=(f,e):(D,s)\to(D',s')$ be an exact morphism and let $Q$ be a core on $(D',s')$.
Writing $F_D(s)$ for the selected extracted family, we have
\begin{equation*}
F_{D'}(s')=e_*F_D(s),\qquad
F_D(s)=(e^{-1})_*F_{D'}(s').
\tag{5.8}\label{eq:5.8}
\end{equation*}
Since the right-hand side is finite, we obtain a finite family on the original as well that
can serve as input to a core. We place the base at $(D,s)$ and define composition, object
formation, and operations as follows.
\begin{equation*}
\begin{aligned}
\mathrm{Comp}_{\sigma^*Q}(F)
 &=(e^{-1})_*\mathrm{Comp}_Q(e_*F),\\
\mathrm{Form}_{\sigma^*Q}(C)
 &=T_{e^{-1}}\mathrm{Form}_Q(e_*C),\\
\mathrm{Op}_{\sigma^*Q}(A,B)
 &=\mathrm{Op}_Q(T_eA,T_eB).
\end{aligned}
\tag{5.9}\label{eq:5.9}
\end{equation*}
Here $T_e$ is the reindexing of objects from Construction~\ref{cons:4.3}. We conjugate the
configuration maps of the operations and compose the invariants and signatures with $T_e$.
The contexts, equations, residuals, and detector code are reindexed by the same $e^{-1}$.
By~\eqref{eq:5.8} and the formulas for composition and object formation, the generated base
object coincides with the base object of $T_{e^{-1}}Q$. This coincidence aligns the base
objects of the equation reading.

The resulting morphism $\kappa_{\sigma,Q}:\sigma^*Q\to Q$ has $\sigma$ as its base. The
upper-level reading map $u$ has a two-sided inverse $u^{-1}$ given by the inverse reindexing.
The information about sources used by this construction is $f(s)=s'$ and the equivalence of
extraction.

\end{construction}

\begin{theorem}[Cartesian lifts along arbitrary exact morphisms]\label{thm:5.6}

For every exact base morphism $\sigma:b\to b'$ and every core $Q$ on $b'$, the morphism
$\kappa_{\sigma,Q}$ of Construction~\ref{cons:5.5} is strongly cartesian with respect to
$q:E_{\mathrm{core}}\to B$.

\end{theorem}

\begin{proof}

For $h:R\to Q$ as in~\eqref{eq:5.7}, let the base of $k$ be $\tau$ and its upper level be
$u^{-1}h_{\mathrm{up}}$. The Atom component agrees with the Atom component of $\tau$, by the fact that
the base of $h$ is $\sigma\tau$ and by $u^{-1}$. The other components are composites of
morphisms of readings, so they preserve object formation, operations, equations, invariants,
and signatures. From $uu^{-1}=1$ we get $\kappa_{\sigma,Q}k=h$. Any other factor also has
base $\tau$, and composing its upper-level equation with $u^{-1}$ gives the same $k$.

\end{proof}

\begin{lemma}[Pullback functors and composition comparisons]\label{lem:5.7}

Choose one cartesian lift for each $\sigma,Q$ by Theorem~\ref{thm:5.6} and fix it as the
canonical choice. For this choice we obtain a functor $\sigma^*:E_{b'}\to E_b$ between fiber
categories. The image of a vertical morphism $v:Q\to Q'$ is determined uniquely by
\begin{equation*}
\kappa_{\sigma,Q'}\,\sigma^*(v)
=v\,\kappa_{\sigma,Q}
\tag{5.10}\label{eq:5.10}
\end{equation*}
For $\sigma:b_0\to b_1$ and $\tau:b_1\to b_2$ there are natural isomorphisms of unit and
composition
\begin{equation*}
1_{E_b}\xrightarrow{\sim}(1_b)^*,
\qquad
\sigma^*\tau^*\xrightarrow{\sim}(\tau\sigma)^*
\tag{5.11}\label{eq:5.11}
\end{equation*}
satisfying the coherence for units and triple composites.

\end{lemma}

\begin{proof}

Applying~\eqref{eq:5.10} to identities and to composites of morphisms and using uniqueness
yields the functor laws. A composite of cartesian morphisms is cartesian, by
applying~\eqref{eq:5.7} twice. Hence between the two-step lift and the direct lift there is a
unique isomorphism making the triangle into their common target commute. The two comparisons
arising from a three-step lift also satisfy the same triangle and are therefore equal. For
the unit we use the same uniqueness.

\end{proof}

\begin{proposition}[Adjoint equivalence between transport and pullback]\label{prop:5.8}

For an exact morphism $\sigma:b\to b'$, the transport of Chapter~\ref{chap:4} and the
pullback of this section form an adjunction
\begin{equation*}
\sigma_!:E_b\rightleftarrows E_{b'}:\sigma^*,
\qquad \sigma_!\dashv\sigma^*
\tag{5.12}\label{eq:5.12}
\end{equation*}
whose unit $\eta^\sigma:1\to\sigma^*\sigma_!$ and counit
$\epsilon^\sigma:\sigma_!\sigma^*\to1$ are invertible.

\end{proposition}

\begin{proof}

Write $\iota_{\sigma,P}:P\to\sigma_!P$ for the opcartesian morphism in the same direction as
in Chapter~\ref{chap:4}. The two universal properties give the following natural bijections.
\begin{equation*}
\begin{aligned}
\mathrm{Hom}_{E_{b'}}(\sigma_!P,Q)
&\cong\{h:P\to Q\mid q(h)=\sigma\}\\
&\cong\mathrm{Hom}_{E_b}(P,\sigma^*Q).
\end{aligned}
\tag{5.13}\label{eq:5.13}
\end{equation*}
The first correspondence is $v\mapsto v\iota_{\sigma,P}$, and the inverse of the second is
$w\mapsto\kappa_{\sigma,Q}w$. The unit and the counit are produced by transposing identity
morphisms. The triangle identities follow from the fact that both composites become identities
after composition with the respective lifts, together with uniqueness.

Moreover, the upper levels of the two lifts are built from mutually inverse Atom reindexings.
Using the proof of Theorem~\ref{thm:5.6} for factorizations in the opposite direction, each
lift is both cartesian and opcartesian. The unit is the comparison of two cartesian lifts
with the same base morphism and the same end, and the counit is the comparison of two
opcartesian lifts with the same base morphism and the same start. Both are invertible by
uniqueness.

\end{proof}

\begin{example}[Fibers become equivalent even when the base is not invertible]\label{ex:5.9}

Take the constant map from the source set $\{0,1\}$ to $\{*\}$ and select $0$ in the former.
Every source extracts the same finite family $F$, and the normalizations and Atom maps are
identities. This is an exact morphism, but it is not an isomorphism of bases, since the
source map has no two-sided inverse. Proposition~\ref{prop:5.8} shows that the fiber
categories over the selected sources are nevertheless equivalent. What travels back and
forth is the object formation, the operations, and the Laws read at that source. Recovering
the distinctions among all sources is a different question.

\end{example}

\section{The canonical comparison exchanging transport and pullback}\label{sec:5.3}

We express the two splitting plans from the opening as routes with the same start and end.
Let $V$ be the route that splits the whole and then selects the scope, and $D$ the route that
selects the scope and then splits; from the universal properties of transport and pullback
we build a comparison between them.

In what follows we assume that Atom equality is decidable. We fix two morphisms presented by
the finite codes of Definition~\ref{def:5.14}, a table of selected sources with a common
image, and a finite diagnostic diagram. The two morphisms have a common target. The square of
this section is the pointed fiber product obtained by decoding these data and applying
Construction~\ref{cons:5.1}. The concrete definition of the finite codes is given
in~\S\ref{sec:5.4}. For transport we use the canonical construction of Chapter~\ref{chap:4},
and for pullback the canonical lifts determined by these finite codes.

\begin{construction}[Two routes with matched types]\label{cons:5.10}

Write the pointed fiber product constructed from these finite codes as
\begin{equation*}
\begin{array}{ccc}
b_{12}&\xrightarrow{\ \pi_2\ }&b_2\\
{\scriptstyle\pi_1}\downarrow&&\downarrow{\scriptstyle\sigma_2}\\
b_1&\xrightarrow{\ \sigma_1\ }&b_0
\end{array}
\qquad \sigma_1\pi_1=\sigma_2\pi_2
\tag{5.14}\label{eq:5.14}
\end{equation*}
Each vertex is a pair of a doctrine and a selected source, and in the opening example the
vertices play the following roles.

\begin{xltabular}{\linewidth}{@{}LL@{}}
\toprule
Vertex & How the input is read \\
\midrule
\endhead
$b_1$ & Reads the input on the monolith side before the split \\
\addlinespace[3pt]
$b_0$ & Reads the whole input after the split \\
\addlinespace[3pt]
$b_2$ & Reads the scope of input selected for order cancellation after the split \\
\addlinespace[3pt]
$b_{12}$ & Reads the pairs of sources, on the monolith side and in the selected scope,
that point to the same target \\
\bottomrule
\end{xltabular}

The morphism $\sigma_1$ is the correspondence given by the splitting policy, and $\sigma_2$
is the correspondence from the selected scope into the whole. The projections $\pi_1,\pi_2$
extract the components from a corresponding pair of sources. The selection of the scope here
is performed on sources, and is distinguished from the operation of removing Atoms inside a
single core. To express these as exact base morphisms, we check the equivalence of
extraction of~\eqref{eq:5.5} for each correspondence.

The two routes from cores on $b_1$ to cores on $b_2$ are
\begin{equation*}
D=(\pi_2)_!\pi_1^*,\qquad
V=\sigma_2^*(\sigma_1)_!,\qquad
D,V:E_{b_1}\longrightarrow E_{b_2}.
\tag{5.15}\label{eq:5.15}
\end{equation*}
The route $D$ pulls back to the corresponding inputs and then carries them over; the route
$V$ carries them to the common base and then pulls back. From the composition comparisons of
transport in Chapter~\ref{chap:4} and~\eqref{eq:5.14} we obtain a natural isomorphism
\begin{equation*}
\Sigma:(\sigma_2)_!(\pi_2)_!
\xrightarrow{\sim}(\sigma_1)_!(\pi_1)_!
\tag{5.16}\label{eq:5.16}
\end{equation*}
It is built by comparing both with the direct transport along the same composite base
morphism.

\end{construction}

\begin{theorem}[Canonical Beck--Chevalley mate]\label{thm:5.11}

The following three-step composite defines a natural isomorphism
$\alpha:D\xrightarrow{\sim}V$.
\begin{equation*}
\begin{aligned}
DP&\xrightarrow{\ \eta^{\sigma_2}_{DP}\ }
 \sigma_2^*(\sigma_2)_!DP\\
&\xrightarrow{\ \sigma_2^*(\Sigma_{\pi_1^*P})\ }
 \sigma_2^*(\sigma_1)_!(\pi_1)_!\pi_1^*P\\
&\xrightarrow{\ \sigma_2^*(\sigma_1)_!(\epsilon^{\pi_1}_P)\ }
 VP.
\end{aligned}
\tag{5.17}\label{eq:5.17}
\end{equation*}
We call this comparison the canonical Beck--Chevalley mate. Moreover, if the choice of
pullbacks is changed, the mates whose two ends are aligned by the unique lift comparisons
coincide.

\end{theorem}

\begin{proof}

The three arrows are, respectively, the images under functors of the unit, the natural
isomorphism of transports, and the counit. Hence the composite is a natural transformation.
By Proposition~\ref{prop:5.8} the unit and the counit are invertible, and functors preserve
isomorphisms, so all three arrows are invertible.

Independence of the choices follows because the next equation characterizes the comparison
itself.
\begin{equation*}
\kappa_{\sigma_2,(\sigma_1)_!P}\,
\alpha_P\,\iota_{\pi_2,\pi_1^*P}
=\iota_{\sigma_1,P}\,\kappa_{\pi_1,P}.
\tag{5.18}\label{eq:5.18}
\end{equation*}
Factoring the right-hand side first through $\iota_{\pi_2,\pi_1^*P}$ and then through
$\kappa_{\sigma_2,(\sigma_1)_!P}$ yields $\alpha_P$ uniquely. Substituting the
factorizations of the unit and counit of the adjunction shows that this morphism
is~\eqref{eq:5.17}. Aligning a different choice of lifts by the canonical isomorphisms
preserves~\eqref{eq:5.18}, so the comparisons coincide.

\end{proof}

The reason for invertibility lies not only in the shape of the fiber product but in the fact
that the unit and counit of the actual transport of cores are invertible. The proof also
shows which components of the comparison have to be checked.

\begin{proposition}[Comparisons under a different choice of lifts]\label{prop:5.12}

In the square of the same finite codes as in Construction~\ref{cons:5.10}, choose different
cartesian lifts along $\pi_1$ and $\sigma_2$. Let $D',V'$ be the resulting two routes and
$\alpha':D'\to V'$ the mate built from the adjunctions. Composing the unique isomorphisms
that connect each lift with the canonical choice, together with the remaining transport
functors, to form $j_D:D'\xrightarrow{\sim}D$ and $j_V:V'\xrightarrow{\sim}V$, we have
\begin{equation*}
\alpha\,j_D=j_V\,\alpha'
\tag{5.19}\label{eq:5.19}
\end{equation*}
\end{proposition}

\begin{proof}

The comparisons of lifts preserve composition with those lifts. Substituting this equation
into the factorizations of the units and counits of the adjunctions, the three-step
composite that builds the modified mate agrees with the composite that places $j_D,j_V$
before and after the canonical mate.

\end{proof}

\begin{example}[A mix-up in implementing the comparison]\label{ex:5.13}

Suppose that, after separating orders and inventory, in comparing the designs obtained by
the two routes the correspondence between two inventory reservations is mixed up. If this
swap is an automorphism $s$ in the chosen reading, set $a_P=s\alpha_P$. Writing the mismatch
of this example as $\Delta_P=a_P\alpha_P^{-1}$, we have $\Delta_P=s$. Here both $a_P$ and
$\alpha_P$ are invertible. Invertibility, and correctly matching the two specified routes,
must each be checked by its own equation.

\end{example}

\section{Existence of the constructions and the scope of finite presentation}\label{sec:5.4}

The pullback of cores in~\S\ref{sec:5.2} exists along every exact morphism. The canonical
comparison of~\S\ref{sec:5.3} used a square generated from finite codes. Here we fix that
format and classify which base morphisms can be presented when the presentations of the
endpoints may be rechosen. The same format is also used for the reconstruction in
Chapter~\ref{chap:8}. The discussion of comparisons and diagnoses continues
in~\S\ref{sec:5.5}.

\begin{definition}[Codes by finite exception tables]\label{def:5.14}

The sources are enumerated as a finite set, and the normalization map is given by a table.
The Atom predicate used at each normalized source is represented by a default truth value
and a finite exception set that flips it. A default value of false represents a finite set,
and a default of true a cofinite set, that is, a set with finite complement. A pointed code
adds the selected source. A morphism between codes uses a finite table of sources and an
Atom permutation table with finite support, and we require commutation with the
normalizations, preservation of the selected sources, and the following equality of codes.
The extraction code obtained by normalizing an original source and transporting it by the
Atom permutation is equal to the extraction code itself obtained by normalizing the
corresponding source of the target. This transport preserves the default value and sends the
finite exception set by the Atom permutation. The equality is therefore required both for
the default values and for the finite exception sets.

An exact morphism $\sigma:b\to b'$ is presentable in this format if there are a morphism
$\widehat\sigma:\widehat b\to\widehat b'$ decoded from a code and isomorphisms of bases
$i:\widehat b\xrightarrow{\sim}b$ and $j:\widehat b'\xrightarrow{\sim}b'$ such that
\begin{equation*}
\sigma i=j\widehat\sigma
\tag{5.20}\label{eq:5.20}
\end{equation*}
The presentation of the whole morphism, the isomorphisms $i,j$ included, is what is being
compared.

\end{definition}

\begin{theorem}[Presentability by this code]\label{thm:5.15}

Assume that Atom equality is decidable. An exact morphism $\sigma=(f,e):b\to b'$ is
presentable in the sense of Definition~\ref{def:5.14} if and only if the following two
conditions hold.

\begin{enumerate}
\item The source sets of both endpoints are finite.
\item For every source $t$ of the target, the extraction set $F_{D'}(t)$ is finite or
cofinite.
\end{enumerate}

\end{theorem}

\begin{proof}

If a presentation exists, the enumeration of the sources and the endpoint isomorphisms give
finiteness of both endpoints. The set represented by a finite exception table is finite or
cofinite, and this property is preserved by bijections of Atoms. The endpoint isomorphisms
match up all sources, so condition~2 holds for every extraction set of the target.

Conversely, assume the two conditions. By exactness, each extraction set of the original is
$e^{-1}_*F_{D'}(f(s))$ and hence also finite or cofinite. We enumerate the sources of both
endpoints and express the normalization and the source map in that enumeration. At sources
in the image of the normalization map, the extraction table encodes the conjunction of the
original predicates. That conjunction equals the extraction predicate of the inputs that
normalize there. The table outside the image may be set to false, because decoding always
normalizes before reading the table. This procedure does not require idempotence of the
normalization.

Choose the Atom components of the endpoint isomorphisms to be the identity on the original
and $e$ on the target. Rereading the tables of the target along $e$ as well, the identity
can be used for the Atom permutation inside the code. For equal extraction sets we choose
the same default value and exception set, making the normalized tables of the original agree
with the corresponding tables of the target. In this way the equality of codes holds in
addition to the equivalence of predicates. The pair with the table of sources is a morphism
in the sense of Definition~\ref{def:5.14} and satisfies~\eqref{eq:5.20}.

\end{proof}

In this conclusion the presentations of the endpoints are also rechosen. Fullness of the Hom
sets---whether, after fixing two codes in advance, every morphism between the decoded
objects can be presented by a morphism of codes---is a further question about those fixed
presentations.

\begin{example}[Finitely many sources are not enough]\label{ex:5.16}

Let the Atom set be $\mathbb{N}$, take a single source, let the normalization be the
identity, and let the extraction set be the even numbers. The identity morphism of this
doctrine is exact, but since both the even and the odd numbers are infinite, it cannot be
presented by a finite exception table. On the other hand, with the same source and
extraction set $\mathbb{N}\setminus\{1,3\}$, it can be presented with default value true and
exceptions $\{1,3\}$. The selected family of the latter is infinite, and no core exists at
that point, since cores of Chapter~\ref{chap:1} require a finite family. The presentation of
the base by finite codes and the existence of a core at the point each have their own
conditions.

\end{example}

\section{Assembling commutative squares into whole diagrams}\label{sec:5.5}

In an actual change we treat not just one square but several APIs, conversions, and
inspection routes at the same time. This requires building correspondences of paths from the
correspondences of the individual edges, and preserving the relations between pairs of
paths.

\begin{definition}[Base diagrams and their changes]\label{def:5.17}

Fix a shape with finitely many vertices, generating edges, and specified pairs of parallel
paths. A base diagram $X$ assigns $X_v\in B$ to each vertex $v$ and $X_a:X_v\to X_w$ to each
edge $a:v\to w$. For the two paths $\ell_f,\rho_f$ of each face $f$ we require
$X_{\ell_f}=X_{\rho_f}$.

A change between diagrams $X,Y$ of the same shape is a family of exact morphisms
$i_v:X_v\to Y_v$, one for each vertex, satisfying for each generating edge
\begin{equation*}
Y_a i_v=i_w X_a
\tag{5.21}\label{eq:5.21}
\end{equation*}
Then $Y_\gamma i_v=i_wX_\gamma$ holds for every path $\gamma:v\to w$ as well. This is proved
by using the identity law for the empty path and substituting~\eqref{eq:5.21} each time the
path is extended by one edge.

What a diagram specifies here is the objects and morphisms of the base and their relations.
Cores, lifts of edges, and automorphisms for comparisons are placed on top of them next.

\end{definition}

\begin{construction}[Building upper-level edges from the same squares]\label{cons:5.18}

Suppose given a core $P_v$ at each vertex and a strongly opcartesian lift $L_a:P_v\to P_w$
for each edge. At the vertices set $P'_v=(i_v)_!P_v$ and write $I_v:P_v\to P'_v$ for the
canonical lift. The edge $L'_a$ on the target is constructed as the unique morphism
satisfying
\begin{equation*}
L'_a I_v=I_wL_a,\qquad q(L'_a)=Y_a
\tag{5.22}\label{eq:5.22}
\end{equation*}
Its existence is given by opcartesianness of $I_v$ and by~\eqref{eq:5.21}. Moreover, since
$I_wL_a$ is opcartesian, $L'_a$ is opcartesian as well. Indeed, compose a morphism out of
the target of $L'_a$ with $I_v$ and factor it through $I_wL_a$. Eliminating $I_v$ from the
resulting equation by the universal property yields the factorization through $L'_a$ and
its uniqueness.

Composing~\eqref{eq:5.22} along a path gives $L'_\gamma I_v=I_wL_\gamma$. For identities,
composites of changes of diagrams, and horizontal and vertical pasting of squares, the two
constructions satisfy the same factorization equations and therefore agree under the
canonical comparisons.

\end{construction}

\begin{proposition}[Conditions for obtaining relations from a family of squares]\label{prop:5.19}

Suppose given only the diagram $X$ on the original, the vertices and edges of the target,
the vertex morphisms $i_v$, and~\eqref{eq:5.21}. The equation obtained for the two paths of
a face $f:v\Rightarrow w$ is
\begin{equation*}
Y_{\ell_f}i_v=i_wX_{\ell_f}
=i_wX_{\rho_f}=Y_{\rho_f}i_v.
\tag{5.23}\label{eq:5.23}
\end{equation*}
If $i_v$ is epi, then $Y_{\ell_f}=Y_{\rho_f}$ follows. Hence, if $i_v$ is epi at every
vertex that appears as the start of a specified face, the family of squares yields a diagram
on the target and a change of diagrams.

Moreover, for a single $i:X\to Y$, the property
\begin{equation*}
ui=vi\Longrightarrow u=v
\tag{5.24}\label{eq:5.24}
\end{equation*}
for all objects $Z$ and all parallel morphisms $u,v:Y\to Z$ is equivalent to $i$ being epi.

\end{proposition}

\begin{proof}

From~\eqref{eq:5.21} for paths and the relations on the original we
obtain~\eqref{eq:5.23}. Applying the cancellation property of epis establishes the relation
for each face. The final equivalence is precisely the definition of an epi.

\end{proof}

The quantification in~\eqref{eq:5.24}, which ranges over all parallel morphisms, differs
from the agreement of only the finitely many paths selected here. For a particular diagram
to be coherent, it is not necessary that each $i_v$ be epi.

\begin{example}[Discrepancy at inputs that are not read]\label{ex:5.20}

Use doctrines that extract the same finite family at every source, with identity
normalizations and Atom maps. Consider the source maps
\begin{equation*}
i:\{*\}\to\{0,1\},\quad i(*)=0,\qquad
u=\mathrm{id}_{\{0,1\}},\quad v(0)=v(1)=0
\tag{5.25}\label{eq:5.25}
\end{equation*}
Taking every selected source to be $*$ or $0$, these become base morphisms. If the two
parallel edges on the original are both identities, the edges $u,v$ on the target
individually form commutative squares. But even though $ui=vi$, we have $u(1)\ne v(1)$.
Agreement at the input $0$ alone does not yield agreement at the input $1$.

If both edges on the target are taken to be $u$, the relation holds even with the same
non-epi $i$. These two examples show the difference between a uniform cancellation
condition and the coherence of an individual diagram.

\end{example}

\section{Carrying diagnoses and edge reselections back and forth}\label{sec:5.6}

Chapter~\ref{chap:4} expressed the difference between a specified comparison and the
canonical comparison as a defect. When the base diagram is changed, we want to carry this
difference as well from the same input. By the fiber equivalence of
Proposition~\ref{prop:5.8}, not only the values of the diagnosis but also the edge
reselections that change those values correspond.

\begin{construction}[End groups, cochains, and the correspondence of reselections]\label{cons:5.21}

Fix the change of diagrams of Definition~\ref{def:5.17} and the cores and edges of
Construction~\ref{cons:5.18}. Let $G_v=\mathrm{Aut}_{E_{X_v}}(P_v)$ be the group of vertical
automorphisms at a vertex $v$, and $G'_v$ its counterpart on the target. The transport
functor induces a group isomorphism
\begin{equation*}
\Phi_v:G_v\xrightarrow{\sim}G'_v,
\qquad \Phi_v(g)=(i_v)_!(g)
\tag{5.26}\label{eq:5.26}
\end{equation*}
Surjectivity comes from fullness and injectivity from faithfulness. A cochain
$\omega=(\omega_f)_f$ on faces and an edge reselection $a=(a_e)_e$ are each sent by applying
this isomorphism at the ends.
\begin{equation*}
(\Phi_C\omega)_f=\Phi_{t(f)}(\omega_f),
\qquad
(\Phi_R a)_e=\Phi_{t(e)}(a_e).
\tag{5.27}\label{eq:5.27}
\end{equation*}
Here $t(f)$ and $t(e)$ are the ends of the face and of the edge. Using each $\Phi_v^{-1}$
componentwise, both assignments have two-sided inverses. The value of a specified comparison
$u_f$ on the target is likewise taken to be $u'_f=\Phi_{t(f)}(u_f)$.

\end{construction}

\begin{theorem}[Preservation and reflection of coherence and reselection orbits]\label{thm:5.22}

Suppose that the base diagram on the target also satisfies the specified path relations. For
the construction above, the canonical comparisons $\phi_f$ of Chapter~\ref{chap:4} and the
raw defects satisfy
\begin{equation*}
\begin{aligned}
\phi'_f(\Phi_Ra)&=\Phi_{t(f)}(\phi_f(a)),\\
\delta'_f(\Phi_Ra)&=\Phi_{t(f)}(\delta_f(a))
\end{aligned}
\tag{5.28}\label{eq:5.28}
\end{equation*}
Consequently, the following equivalences hold.

\begin{enumerate}
\item All faces are coherent at $a$ if and only if all faces are coherent at $\Phi_Ra$.
\item A reselection making all faces coherent exists before the change if and only if one
exists after the change.
\item $\omega$ belongs to the reselection orbit of the raw defect if and only if
$\Phi_C\omega$ belongs to the orbit after the change.
\end{enumerate}

For an arbitrary cochain after the change as well, the same verdicts hold for the value
brought back by $\Phi_C^{-1}$.

\end{theorem}

\begin{proof}

A reselected edge is $L_e^a=a_eL_e$. By the action of transport on morphisms
and~\eqref{eq:5.22}, its image is $\Phi_{t(e)}(a_e)L'_e$. Extending to paths, the morphism
obtained by carrying the canonical comparison of the original satisfies the equation comparing the
two paths after the change. By uniqueness of strongly opcartesian lifts, this is the
canonical comparison after the change. Since group homomorphisms preserve products and
inverses, the second equation of~\eqref{eq:5.28} follows from $\delta_f=u_f\phi_f^{-1}$.

Each $\Phi_v$ preserves and reflects the identity element, so the defect being the identity
in all components is equivalent on the two sides. Moreover $\Phi_R$ is a bijection, so every
reselection after the change can be brought back as well. Carrying the witness
$\omega=\delta(a)$ of orbit membership by~\eqref{eq:5.28}, and using $\Phi_R^{-1}$ and
$\Phi_C^{-1}$ in the reverse direction, the three conclusions follow.

\end{proof}

This diagnosis concerns the fixed diagram and the comparisons generated from it. When new
faces or new Laws are added after the change, a diagram including them is specified anew.
For geometries involving a coefficient ring, this section keeps the same coefficient ring
fixed.

\begin{example}[Correspondence of reselections that cancel the defect]\label{ex:5.23}

In the one-object group category example of Chapter~\ref{chap:4}, take the automorphism
group to be the permutation group $S_3$ on three letters. On two parallel identity edges,
place the specified comparison $u=(12)$. Reselecting the left edge to $a=1$ and the right
edge to $b=(12)$, the canonical comparison becomes $ba^{-1}=(12)$ and the defect vanishes.

Under the group isomorphism $\Phi(g)=cgc^{-1}$ given by the renaming $c=(123)$, the
specified comparison and the reselection of the right edge are both carried to $(23)$. Hence
the same face is coherent after the change. Not only is the truth value ``vanished''
preserved; the reselection that achieved it is carried along.

\end{example}

\section{Conditions for backward transport along refinements}\label{sec:5.7}

When the extraction rules are made more detailed, Atoms that were not extracted before may
appear. Such a change cannot always be carried backwards in the way an exact morphism can.
What is needed is that, at the sources where cores are actually placed, the extraction on
the target can be brought back to the original.

\begin{definition}[Refinements and the realized locus]\label{def:5.24}

A refinement of bases $r:(D,s)\to(D',s')$ has a source map $f$ and a bijection $e$ of
Atoms, and satisfies $f(s)=s'$, $fN_D=N_{D'}f$, and the forward implication
\begin{equation*}
\mathrm{Extracts}_D(t,a)
\Longrightarrow\mathrm{Extracts}_{D'}(f(t),e(a))
\tag{5.29}\label{eq:5.29}
\end{equation*}
Exactness is the condition that adds the backward implication to this.

The collection of objects of the base that carry at least one core is called the realized
locus. We say that $r$ reflects extraction on the realized locus if, whenever $E_{b'}$ has
an object, for every Atom $a$ we have
\begin{equation*}
\mathrm{Extracts}_{D'}(s',e(a))
\Longrightarrow\mathrm{Extracts}_D(s,a)
\tag{5.30}\label{eq:5.30}
\end{equation*}
We write $R(r)$ for this conditional property, namely ``if $E_{b'}$ has an object,
then~\eqref{eq:5.30} holds''.

For core morphisms over refinements, only the base is relaxed to~\eqref{eq:5.29}, while the
upper level is subject to the same preservation conditions for readings as in
Definition~\ref{def:1.28}. In particular, the comparison of selected families is still
required as an equality. These morphisms and objects are closed under composition and
define a category $E_{\mathrm{core}}^{\mathrm{ref}}$ with a projection $q_{\mathrm{ref}}$
to the base.

\end{definition}

\begin{theorem}[Classification of cartesian lifts by the realized locus]\label{thm:5.25}

Fix one pointed refinement $r:b\to b'$. That a strongly cartesian lift with respect to
$q_{\mathrm{ref}}$ can be chosen for every core on $b'$ is equivalent to $R(r)$.

\end{theorem}

\begin{proof}

Assume $R(r)$ and take a core $Q$ on $b'$. Combining~\eqref{eq:5.29} and \eqref{eq:5.30} gives
$F_{D'}(s')=e_*F_D(s)$. From this equality we bring finiteness back and build a core $P$ by
the same inverse reindexing as in Construction~\ref{cons:5.5}. The upper level of the
morphism $P\to Q$ with base $r$ has a two-sided inverse. The same factorization as in
Theorem~\ref{thm:5.6} works after an arbitrary refinement base morphism as well.

Conversely, suppose that lifts can be chosen, and take one core $Q$ on $b'$. The upper
level of its lift $P\to Q$ satisfies $F_Q=e_*F_P$ for the selected families. If $e(a)$ is
extracted on the target, then $a\in F_P$ by injectivity of $e$, so it is extracted on the
original as well. This is~\eqref{eq:5.30}. When there is no core on $b'$, there is no
object for which to choose a lift, and both conditions hold vacuously.

\end{proof}

\begin{construction}[Forward base change of refinements]\label{cons:5.26}

Take an exact cospan $D_1\to D_0\leftarrow D_2$ and a refinement $r:D'_1\to D_1$. Define
the pairs of corresponding sources by
\begin{equation*}
\mathrm{Src}'_{12}
=\{(s'_1,s_2)\mid f_1(f_r(s'_1))=f_2(s_2)\}
\tag{5.31}\label{eq:5.31}
\end{equation*}
and define extraction from the first component in $D'_1$. The first projection
$\pi'_1:D'_{12}\to D'_1$ is exact, and $(s'_1,s_2)\mapsto(f_r(s'_1),s_2)$ together with the
Atom map $e_r$ defines a refinement $r_{12}:D'_{12}\to D_{12}$. Preservation of extraction
follows from~\eqref{eq:5.29}, and the following square commutes.
\begin{equation*}
\begin{array}{ccc}
D'_{12}&\xrightarrow{\ r_{12}\ }&D_{12}\\
{\scriptstyle\pi'_1}\downarrow&&\downarrow{\scriptstyle\pi_1}\\
D'_1&\xrightarrow{\ r\ }&D_1.
\end{array}
\tag{5.32}\label{eq:5.32}
\end{equation*}
This construction does not require $R(r)$. Choosing one pair in~\eqref{eq:5.31} fixes the
sources of the four vertices and a pointed square.

\end{construction}

\begin{proposition}[Transfer of the realized locus and the two backward routes]\label{prop:5.27}

Choose a compatible pair in~\eqref{eq:5.31} and suppose that $R(r)$ holds at that point.
Then $R(r_{12})$ holds as well, and we obtain two functors from cores on $D_1$ to cores on
$D'_{12}$ and a natural comparison
\begin{equation*}
(\pi'_1)^*r^*
\longrightarrow r_{12}^*\pi_1^*
\tag{5.33}\label{eq:5.33}
\end{equation*}
That both refinement lifts can be chosen at every compatible pair is equivalent to $R(r)$
holding at all such pairs. Moreover, pointed refinements that reflect extraction on the
realized locus are closed under identities and composition.

\end{proposition}

\begin{proof}

If there is a core at the selected point of $D_{12}$, it can be transported forwards along
the exact $\pi_1$. This yields a core at the corresponding point of $D_1$ as well, and
$R(r)$ can be used. The extraction of both fiber products is that of the first component,
so~\eqref{eq:5.30} carries over to $r_{12}$ as it stands. Theorem~\ref{thm:5.25} gives the
two refinement lifts, and Theorem~\ref{thm:5.6} the two exact lifts. Factoring the two-step
lifts along the commutative square gives~\eqref{eq:5.33}. Its naturality follows because
the two factors obtained by composing with a vertical morphism satisfy the same triangle.
For the equivalence over all pairs we apply Theorem~\ref{thm:5.25} pointwise.

For a composite $b_0\xrightarrow{r}b_1\xrightarrow{t}b_2$, a core on $b_1$ is built by
pulling a core on $b_2$ back along $t$. This also satisfies the realization condition of
$R(r)$. Reflecting extraction along $t$ and then along $r$ yields $R(tr)$. In the case of
an identity, \eqref{eq:5.30} becomes the identity implication.

\end{proof}

\begin{example}[A forward-only change and the scope that can be carried back]\label{ex:5.28}

Let the Atoms be $\{A,B,C\}$ and the sources $\{\mathrm{partial},\mathrm{all}\}$. Take the
identity normalization and define the extractions of the original and the target as
follows.

\begin{xltabular}{\linewidth}{@{}LLL@{}}
\toprule
Source & Original & Target \\
\midrule
\endhead
partial & A, B & A, B, C \\
\addlinespace[3pt]
all & A, B, C & A, B, C \\
\bottomrule
\end{xltabular}

The source map is the identity, and the Atom map is the bijection that swaps $A,B$ and
fixes $C$. This is a refinement. All extracted families are finite, and a core can be
placed at each point. For instance, set $\mathrm{Comp}(F)=(F,\varnothing,\varnothing)$, and
let object formation preserve the configuration and place the element of a one-element set
in the structure data and in the quantity. If the operations are only the identities
between equal objects, and the equation indices, invariants, and signature axes are empty,
we can construct a core satisfying the conditions of the reading.

If partial is selected, the Atom $C$ on the target cannot be brought back to the
original. By Theorem~\ref{thm:5.25}, no lift to a core on the target exists. If both cospan
legs are the identity of the target doctrine, the forward square of
Construction~\ref{cons:5.26} can still be built.

On the other hand, let the second cospan leg be the exact morphism into all on the target
from the doctrine that has a single source and extracts all Atoms. In a compatible pair
the first source is then restricted to all, where~\eqref{eq:5.30} holds. Hence that square
has both backward transports. The original refinement still enlarges the extraction at
partial. The reason backward transport became possible is that the inputs matched this
time are narrowed to the scope of all.

\end{example}

\begin{example}[Contrast with the case where no object exists]\label{ex:5.29}

Consider the identity refinement with one source, Atoms $\mathbb{N}$, and the whole of
$\mathbb{N}$ as the extracted family. Since the selected family is infinite, no core
exists. Both conditions of Theorem~\ref{thm:5.25} hold, but this is not an example of cores
actually traveling back and forth. Example~\ref{ex:5.28} compares the two behaviors after
confirming that objects exist.

\end{example}

\section{Comparisons including covers and local data}\label{sec:5.8}

A comparison of cores includes the objects, the operations, and the Laws. To compare local
inspections under the same conditions, we must also carry which contexts read what and
where they overlap. The full geometries of Chapter~\ref{chap:1} are the objects that add
these covers, overlaps, coefficients, and raw systems to a core.

\begin{definition}[Morphisms of geometries over refinements]\label{def:5.30}

For a morphism of geometries $G\to H$, take its lower level to be a morphism of
$E_{\mathrm{core}}^{\mathrm{ref}}$, and take the correspondences of contexts, equations,
Atoms, and axes from its upper-level reading map. On the components of the geometries we
impose the same five conditions as in Definition~\ref{def:1.30}: forward preservation of
coverage requirements, comparison of overlaps, a coefficient homomorphism, the transport
equations for raw systems, and the readouts of support, axes, and observables together with
their naturality. These define a category $E_{\mathrm{geom}}^{\mathrm{ref}}$ and
projections
\begin{equation*}
E_{\mathrm{geom}}^{\mathrm{ref}}
\xrightarrow{\ p_{\mathrm{ref}}\ }
E_{\mathrm{core}}^{\mathrm{ref}}
\xrightarrow{\ q_{\mathrm{ref}}\ }B^{\mathrm{ref}}
\tag{5.34}\label{eq:5.34}
\end{equation*}
The exact three-level categories enter here by restricting the base to exact morphisms.

\end{definition}

\begin{construction}[Pullback of geometries along the generated backward transport]\label{cons:5.31}

Take a lift of cores $K:P\to Q$ built by Theorem~\ref{thm:5.6} or 5.25, and a geometry $H$
on $Q$. Let $u,u^{-1}$ be the upper level of $K$ and its generated inverse, and
$\theta,\bar\theta$ the correspondence of contexts. We build a geometry $K^*H$ on $P$ as
follows.

\begin{xltabular}{\linewidth}{@{}LL@{}}
\toprule
Data & How it is pulled back \\
\midrule
\endhead
Coverage requirements and visibility & Send Atoms, equations, and axes forward, then read
the predicates of $H$ \\
\addlinespace[3pt]
Overlap of two contexts & Send them to $H$ by $\theta$, take the overlap in $H$, and return
by the specified quasi-inverse \\
\addlinespace[3pt]
Coefficient ring & Use the same ring $k$ as $H$ \\
\addlinespace[3pt]
Raw system & In a context $W$, read the coordinates, relations, and local data of $H$ at
$\theta W$ \\
\addlinespace[3pt]
Support, axes, observables & Use the same correspondence of carriers that accompanies the
reindexing \\
\bottomrule
\end{xltabular}

For the restriction maps of the raw system we use the morphisms sent by $\theta$. The
identity and composition laws and the preservation of relation polynomials follow from the
original raw system. The universal property of overlaps is brought back by the equivalence
of context categories. In this way we obtain a morphism with identity coefficient
component,
\begin{equation*}
\widehat K:K^*H\longrightarrow H
\tag{5.35}\label{eq:5.35}
\end{equation*}
This morphism is strongly cartesian with respect to $p_{\mathrm{ref}}$. Indeed, for a
morphism passing through $K$ on the lower level, the cover conditions return to the
preimage conditions just defined, and the raw-system equations return by the reindexing of
$u^{-1}$. The comparison maps of local data are also brought back by the generated
two-sided inverses. These constitute the factor, and uniqueness follows by composing with
the same inverses.

\end{construction}

\begin{theorem}[Two backward routes built from the same geometry]\label{thm:5.32}

Fix the point and the realized-locus condition of Proposition~\ref{prop:5.27}, and place a
geometry $H$ on a core $Q$ on $D_1$. Repeating Construction~\ref{cons:5.31} along the two
routes yields geometries $H_B,H_P$ and morphisms $K_B:H_B\to H$ and $K_P:H_P\to H$. Between
them there is a natural comparison $m_H:H_B\to H_P$ satisfying
\begin{equation*}
K_Pm_H=K_B,
\qquad p_{\mathrm{ref}}(m_H)=m_Q
\tag{5.36}\label{eq:5.36}
\end{equation*}
Here $m_Q$ is the comparison of cores obtained by aligning the endpoints of the two routes
of~\eqref{eq:5.33} by the canonical isomorphisms. For the generated routes $m_H$ is
invertible, and its coefficient component is $1_k$.

\end{theorem}

\begin{proof}

The core morphisms of the two routes are each composites of two strongly cartesian
morphisms. Since the square of bases commutes, the comparison $m_Q$ aligning their starts
and ends, and its inverse, are obtained uniquely. The two routes of geometries are likewise
composites of cartesian morphisms from Construction~\ref{cons:5.31}. We first build $m_Q$
on the lower level and lift its factorization to the upper level to obtain $m_H$.
Performing the same procedure in the opposite direction and applying uniqueness gives a
two-sided inverse. The coefficient map at each step is the identity, so the coefficient map
of the comparison is the identity as well.

\end{proof}

To compare a whole diagram, we take a finite rooted diagram with a path to each vertex, and
fix a diagram of cores inside the fiber category at the same point of the base, geometries
on it, and a coefficient ring $k$. We further give the two-level transport data of
Chapter~\ref{chap:4}: each edge of cores and of geometries is a strongly opcartesian
morphism with respect to its projection, and the coefficient components of the edges and of
the specified comparisons are identities. Applying Theorem~\ref{thm:5.32} at each vertex,
we pull the edges and the specified comparisons of the same source back to both routes. As
a result, the generated comparison matches up not only the edges but also the specified
comparisons of the faces.
\begin{equation*}
\begin{aligned}
m_w L^B_e&=L^P_e m_v,\\
m_{t(f)}u^B_f&=u^P_fm_{t(f)}.
\end{aligned}
\tag{5.37}\label{eq:5.37}
\end{equation*}
Composed with $K_P$, both sides of the second equation become the same morphism, obtained by carrying
a single source comparison. The equation is obtained by reflecting it first through
cartesianness of cores and then through cartesianness of geometries. The first equation is
similar, and naturality for paths follows by induction from the generating edges. When the
specified comparisons are chosen separately on the two routes, the second equation becomes
an additional condition to be checked.

\begin{lemma}[Transfer of solutions by endpoint isomorphisms]\label{lem:5.33}

Suppose that between the diagrams of the two routes and other presentations of them there
are natural isomorphisms
$b_v:\widetilde H_{B,v}\xrightarrow{\sim}H_{B,v}$ and
$c_v:\widetilde H_{P,v}\xrightarrow{\sim}H_{P,v}$.
Assume the isomorphisms fix the coefficients and match up the morphisms of the routes and
the specified comparisons. The families of comparisons $s_v$ satisfying the fixed
lower-level comparison, the triangles, and~\eqref{eq:5.37} then correspond in both
directions by
\begin{equation*}
\widetilde s_v=c_v^{-1}s_vb_v,
\qquad
s_v=c_v\widetilde s_vb_v^{-1}
\tag{5.38}\label{eq:5.38}
\end{equation*}
The endpoints of the lower-level comparisons and of the triangles are aligned by the same
isomorphisms.

\end{lemma}

\begin{proof}

Substitute the naturality of $b,c$ into the naturality for edges and cancel adjacent
inverses. The same computation applies to the equations for the specified comparisons and
to the triangles. Using the identity law for the empty path and associativity for
concatenation of paths and pasting of faces, all conditions are preserved. The two formulas
of~\eqref{eq:5.38} are mutually inverse.

\end{proof}

Using the same conjugation at the ends for reselections and defects, vanishing and
reselection orbits correspond in both directions, for the same reason as
in~\eqref{eq:5.28}. What this correspondence uses is routes generated from the same source
and endpoint isomorphisms that preserve the components. For an arbitrarily given morphism
of geometries, the preservation of coverage requirements and local data may be forward
only, and transferring solutions backwards requires the inverse morphisms that
form~\eqref{eq:5.38}.

\section{Reading the exchange of order in lenses and protocols}\label{sec:5.9}

When the view of a lens or the operations of a protocol are restricted, we ask whether
semantics-preserving changes can also be restricted to that scope. For lenses we show that
after the restriction the updates stay within the scope; for protocols, that the
comparisons of executions and adapters are preserved.

\begin{proposition}[Restriction of the view of a lens]\label{prop:5.34}

For a total lens $g:C\to V$, $p:C\times V\to C$ satisfying the three laws and a subset
$W\subseteq V$, set
\begin{equation*}
C_W=\{c\in C\mid g(c)\in W\}
\tag{5.39}\label{eq:5.39}
\end{equation*}
The restrictions of read and update define a total lens $C_W\to W$.

\end{proposition}

\begin{proof}

The read after an update is $g(p(c,w))=w$, so if $w\in W$ the state after the update also
belongs to $C_W$. Reading the equations of the three laws on this subset gives the three
laws after the restriction.

\end{proof}

\begin{example}[Restriction and update with four states]\label{ex:5.35}

We compute the restriction of views and the correspondence of states on a product lens
with four states. Let the displayed values be $V=\{0,1\}$ and the component preserved by
updates be $R=\{0,1\}$, and use the product lens $C=V\times R$. Read returns the first
component, and update rewrites only the first component. Let the displayed values of the
target be $V'=\{A,B\}$, with the bijection $u$ sending 0 to A and 1 to B. The
correspondence of states $h$ converts the displayed value by $u$ and swaps 0 and 1 in the
second component. Let $W'$ be the set of displayed values to select, keeping only A on the
target. That is, we set
\begin{equation*}
h(v,r)=(u(v),1-r),\qquad W'=\{A\}
\tag{5.40}\label{eq:5.40}
\end{equation*}
Selecting the states with displayed value 0 on the original and then applying the
correspondence, or selecting the states with displayed value A on the target, reaches the
same two states. The restricted update is closed in the selected view, and $h$ preserves
updates as well. In this example, comparing the two routes specifies at the same time
which states to keep and which updates to allow.

\end{example}

\begin{proposition}[Restriction of protocols and adapters]\label{prop:5.36}

Let $\mathcal{P}$ be the presented protocol category of Chapter~\ref{chap:1}, and let
$j:\mathcal{Q}\to\mathcal{P}$ be the inclusion of the protocol consisting of the vertices
and operations to be kept and the relations between them. Two realizations
$F,F':\mathcal{P}\to{\mathbf{Set}}$ and a semantics-preserving adapter $a:F\Rightarrow F'$
determine restrictions $Fj,F'j$ and an adapter $aj:Fj\Rightarrow F'j$. Comparison triangles
and commutative adapter squares are also preserved by the restriction.

\end{proposition}

\begin{proof}

The restrictions $Fj,F'j$ are given by composition of functors. The naturality of $a$ for
an arbitrary morphism $e:x\to y$ of $\mathcal{Q}$, namely
$a_{j(y)}F(j(e))=F'(j(e))a_{j(x)}$, gives the naturality of $aj$. Preservation of
observations and the equations of the triangles and squares are likewise restricted at
each vertex.

\end{proof}

In the example of splitting the monolith, from the protocol of the whole we keep the
cancellation request, the release of inventory reservations, and cancellation completion,
together with the relations between them. Giving realizations before and after the split
and checking the naturality of the adapter for each operation, we obtain a correspondence
for the execution of a cancellation that chains them. In selecting this part, we include
the operations carried by the inventory service. Only then does the cancellation process
that crosses services become an object of comparison. The Laws on the states of
cancellation completion and reservation release are checked together on the corresponding
states.

When communication failures and retries are read as new operations and states, we specify
a protocol and a correspondence that include them. Coherence of the comparison is a
property of the selected operations and relations, and for diagrams that also contain
operations outside the selection, it must be checked separately, as in
Example~\ref{ex:5.20}.

\section{The generated comparison including the two routes and normalization}\label{sec:5.10}

The canonical comparison $\alpha$ gives the correspondence of structure along the two
routes. The generated comparison $\beta$ additionally incorporates a normalization
selected according to the diagnosis. Even when a normalization preserves the operations and
the Laws, the distinctions among the original representations may not be recoverable. We
construct this comparison from the same input, and Chapter~\ref{chap:6} analyzes its
invertibility and the information that remains.

\begin{construction}[Two routes of full geometries]\label{cons:5.37}

In this section we fix the following input.

\begin{itemize}
\item Decidable Atom equality, and the square~\eqref{eq:5.14} obtained by decoding the
finite codes of~\S\ref{sec:5.3}.
\item A core at each vertex of the finite diagnostic diagram included in those codes, and a
strongly opcartesian lift for each edge, where the strength is with respect to the projection to
the extraction base.
\item That the images in the base of the two paths of each face are equal, and that the
core at the end of each face lies over $b_1$.
\item A specified automorphism at the end of each face, fixing the extraction base.
\end{itemize}

Choose a face $z$ of these same data and write $P_z$ for the core at its end. Fix a
commutative coefficient ring $k$ and a full geometry $G_z$ on $P_z$. The geometry includes
the coverage requirements, the overlaps, and the raw system. The cochain $\omega$ used
later is also taken for this same diagram and these cores. Using the canonical transport of
geometries of Chapter~\ref{chap:4} and Construction~\ref{cons:5.31}, we form
\begin{equation*}
\overline D(G_z)=(\pi_2)_!\pi_1^*G_z,
\qquad
\overline V(G_z)=\sigma_2^*(\sigma_1)_!G_z
\tag{5.41}\label{eq:5.41}
\end{equation*}
The lower and upper indices denote the operations on entire geometries. Both routes keep
the coefficient ring $k$.

We construct the canonical comparison of full geometries from the comparison of the two
backward routes of~\S\ref{sec:5.8}. Setting $H_0=(\sigma_1)_!G_z$, the comparison that
traces the same square backwards is
\[
m_{H_0}:\pi_1^*\sigma_1^*H_0
\xrightarrow{\sim}\pi_2^*\sigma_2^*H_0
\]
Aligning the choices at the endpoints, the following three steps can be composed using the
unit and the counit.
\[
\begin{aligned}
\overline D(G_z)
&\xrightarrow{\ (\pi_2)_!\pi_1^*(\eta^{\sigma_1}_{G_z})\ }
(\pi_2)_!\pi_1^*\sigma_1^*H_0\\
&\xrightarrow{\ (\pi_2)_!(m_{H_0})\ }
(\pi_2)_!\pi_2^*\sigma_2^*H_0\\
&\xrightarrow{\ \epsilon^{\pi_2}_{\sigma_2^*H_0}\ }
\overline V(G_z).
\end{aligned}
\]
The first morphism is the unit pulled back and then transported, and the last is the
counit. When a different choice of lifts is used, the endpoint isomorphisms built from
Construction~\ref{cons:5.31} are inserted before and after the middle comparison. The order
is: unit, isomorphism at the start, comparison of the backward routes, isomorphism at the
end, counit. Substituting~\eqref{eq:5.36} and the unit-counit triangles yields the
geometric version of~\eqref{eq:5.18}, and we write $\overline\alpha_z$ for the resulting
isomorphism.

\end{construction}

\begin{proposition}[Compatibility with the projection to cores]\label{prop:5.38}

Write $p$ for the functor that forgets the geometry, and
$j_D:p\overline D(G_z)\xrightarrow{\sim}DP_z$ and
$j_V:p\overline V(G_z)\xrightarrow{\sim}VP_z$ for the generated endpoint isomorphisms.
Then
\begin{equation*}
j_V\,p(\overline\alpha_z)=\alpha_{P_z}\,j_D.
\tag{5.42}\label{eq:5.42}
\end{equation*}

\end{proposition}

\begin{proof}

Projecting the lifts of geometries yields the lifts of cores used in the generation. Even
when the choices differ, $j_D,j_V$ are their unique comparisons. Projecting the geometric
version of~\eqref{eq:5.18} and aligning it by these comparisons gives~\eqref{eq:5.18} for
cores. That morphism is unique, so it coincides with the mate of Theorem~\ref{thm:5.11}.

\end{proof}

So far, the structures obtained by the two routes have been matched by an invertible
canonical comparison. Next we choose which differences of representation to keep in the
comparison. For example, we may want to align several design representations with the same
configuration after the split to a common standard representation. To add this processing,
we must preserve not only the presentation of objects but also the operations and the
Laws. We first extract the conditions under which such a normalization can be defined.

\begin{definition}[Normalization preserving configurations]\label{def:5.39}

A core $P$ selects an object $\mathrm{Form}_P(C)$ for each configuration $C$. Let the map
sending an arbitrary object $A$ to the selected object on the same configuration be
\begin{equation*}
n_P(A)=\mathrm{Form}_P(C_A)
\tag{5.43}\label{eq:5.43}
\end{equation*}
By the conditions on object formation, $C_{n_P(A)}=C_A$ and $n_P^2(A)=n_P(A)$. To make
this map a morphism of cores, we let $\mathrm{Ad}(P)$ be the following conditions.

\begin{itemize}
\item In every context, object, equation index, and Atom, the residual is unchanged under
$A\mapsto n_P(A)$.
\item $\mathrm{Op}_P(A,B)$ and $\mathrm{Op}_P(n_P(A),n_P(B))$ have the same type, and this
identification preserves the configuration maps.
\item The indices and kinds of the invariants are preserved, and on every object the truth
values of predicate-type invariants and the values of function-type invariants are
unchanged by the normalization.
\item On every signature axis, the coordinate values are unchanged under
$A\mapsto n_P(A)$.
\end{itemize}

The function-type condition is equivalent to the correspondence by bijections of the value
ranges in Definition~\ref{def:1.28}. Indeed, abbreviate $n=n_P$ and suppose there are
bijections $t_i$ with $t_i(I_i(A))=I_i(n(A))$ for all objects. By idempotence $n^2=n$ of
the object map,
\[
t_i(I_i(A))=I_i(n(A))=I_i(n(n(A)))=t_i(I_i(n(A))).
\]
Injectivity of $t_i$ gives $I_i(A)=I_i(n(A))$. Conversely, if the values are unchanged, we
may choose $t_i=\mathrm{id}$.

We then define an endomorphism $N_P$ of the core by using $n_P$ on objects, the type
identification above on operations, and identities on Atoms, contexts, equation indices,
and axes. Since configurations are preserved, the detector code is also kept identical.
For a geometry $G$ we further define an endomorphism $N_G$ using the identity maps on
coefficients, support, axes, and observables. The equations for covers, overlaps, and raw
systems hold with identity comparisons.

This normalization is incorporated into the comparison of the two routes according to a
specified diagnosis. In the next construction, the normalization is selected when a
component of the cochain is not the identity and the preservation conditions above hold.
In this way, which diagnosis was responded to, and on which route the representations were
aligned, can be expressed as a single comparison morphism.

\end{definition}

\begin{construction}[Generated comparisons selected by the same diagnosis]\label{cons:5.40}

We use the common input of Construction~\ref{cons:5.37} and a cochain $\omega$ on the same
diagram. For the chosen face $z$ and its end $P_z,G_z$, the endomorphism used on the
original core is selected by
\begin{equation*}
S_{z,\omega}=
\begin{cases}
N_{P_z},&\omega_z\ne1\ \text{and}\ \mathrm{Ad}(P_z),\\
1_{P_z},&\text{otherwise}
\end{cases}
\tag{5.44}\label{eq:5.44}
\end{equation*}
This selection is a mathematical case distinction and does not assume a finite decision
procedure for $\mathrm{Ad}(P_z)$. For the full geometry we likewise select $N_{G_z}$ or
the identity under the same condition and write $\overline S_{z,\omega}$. Next, this
endomorphism is carried along the actual second route.
\begin{equation*}
\begin{aligned}
E_{z,\omega}&=V(S_{z,\omega}):VP_z\longrightarrow VP_z,\\
\overline d_{z,\omega}
 &=\overline V(\overline S_{z,\omega}):
 \overline V(G_z)\longrightarrow\overline V(G_z),\\
\beta_{z,\omega}&=E_{z,\omega}\alpha_{P_z},\\
\overline\beta_{z,\omega}
 &=\overline d_{z,\omega}\overline\alpha_z.
\end{aligned}
\tag{5.45}\label{eq:5.45}
\end{equation*}
This is one generation rule that incorporates the permitted normalization when the
component of the specified cochain is not the identity. The current raw defect may also be
chosen as the cochain. The square, the face, the cochain, the coefficients, and the
original geometry used in the construction are common to the two routes. Under the same
projections and endpoint isomorphisms,
\begin{equation*}
j_V\,p(\overline d_{z,\omega})=E_{z,\omega}j_V,
\qquad
j_V\,p(\overline\beta_{z,\omega})=\beta_{z,\omega}j_D
\tag{5.46}\label{eq:5.46}
\end{equation*}
The first equation follows from having selected the same normalization on the original
object and from the projection compatibility of transport. The second is the composite of
the first with~\eqref{eq:5.42}.

Three kinds of comparisons are now in place.

\begin{xltabular}{\linewidth}{@{}LLL@{}}
\toprule
Comparison & What determines it & What to check \\
\midrule
\endhead
The canonical $\alpha$ & The universal properties of transport and pullback & It satisfies
the triangle of the two routes and is invertible \\
\addlinespace[3pt]
A separately specified $a$ & A comparison chosen by an implementation or a way of reading &
Whether it satisfies the triangle, and what its difference from $\alpha$ is \\
\addlinespace[3pt]
The generated comparison $\beta$ & The $\alpha$ of the same input, and the normalization
selected by the diagnosis & Invertibility and image after the normalization is
incorporated \\
\bottomrule
\end{xltabular}

Chapter~\ref{chap:6} proves the idempotence of $N_P$ and $N_G$ as morphisms and studies
when this generated comparison is invertible and, when it is not, what it leaves as its
image.

\end{construction}

\section*{Summary of the Chapter}

In this chapter we built two routes from transport and changes of input, and extended
their comparison to diagnoses, geometries, and normalization.

\begin{itemize}
\item \textbf{Base change and the exchange of order.} We built the inputs corresponding to
a common condition as fiber products and pulled structure back. Along an exact base
morphism a cartesian lift to any core on the target is obtained, and transport and
pullback form an adjoint equivalence. The Beck--Chevalley mate built from its unit and
counit is invertible (\S\S\ref{sec:5.1}--\ref{sec:5.3}).
\item \textbf{Correspondence of diagrams and diagnoses.} In diagrams whose base relations
are aligned, the end groups, cochains, and edge reselections correspond, preserving and
reflecting the vanishing of defects and membership in reselection orbits. In aligning the
relations of a whole diagram, we distinguish conditions cancellable by epis from relations
imposed individually (\S\S\ref{sec:5.5}--\ref{sec:5.6}).
\item \textbf{The scope of backward transport.} For refinements, reflection of extraction
on the realized locus is necessary and sufficient for backward transport. Even when the
doctrine as a whole is not exact, backward transport is possible if the extractions agree
at the sources actually matched (\S\ref{sec:5.7}).
\item \textbf{Extension to geometries and operations.} By carrying the covers and local
data as well from the same original geometry, we raised the comparison of the two routes
to full geometries. For lenses we showed the condition that updates stay within the
selected scope, and for protocols the correspondence of all executions of the kept
operations (\S\S\ref{sec:5.8}--\ref{sec:5.9}).
\item \textbf{Generated comparisons including normalization.} We selected a normalization
according to the diagnosis and composed it with the canonical comparison to form
$\beta=E\alpha$. Since the same normalization is selected for the geometry and the core,
this comparison is also compatible with the projection (\S\ref{sec:5.10}).
\end{itemize}

For the split of the monolith from the opening, we compare the splitting proposal built
from the order-cancellation part with the cancellation part extracted after splitting the
whole. Including the cancellation request, the release of inventory reservations, and
cancellation completion in the scope makes the operations and Laws of the two routes
correspond. If, on top of this, several design representations are aligned to a common
representative, the generated comparison provides the correspondence including that
normalization.

The next chapter examines the normalization factor of the generated comparison and
clarifies its invertibility and the information that remains.
